% Translated front matter. Source PDF pages 1--15.
\pagenumbering{roman}

\begin{titlepage}
\centering
\vspace*{2cm}
{\Large 应用数学教材丛书 15\par}
\vspace{2cm}
{\huge\bfseries 有限元方法的数学理论\par}
\vspace{0.8cm}
{\Large 第三版\par}
\vfill
{\Large Susanne C. Brenner\qquad L. Ridgway Scott\par}
\vspace{1.5cm}
{\Large Springer\par}
\end{titlepage}

\clearpage
\thispagestyle{empty}
\begin{center}
{\Large\bfseries 应用数学教材丛书 15}\par
\vspace{1.5cm}
{\bfseries 编辑}\par
J.E. Marsden\par
L. Sirovich\par
S.S. Antman\par
\vspace{1cm}
{\bfseries 顾问}\par
G. Iooss\par
P. Holmes\par
D. Barkley\par
M. Dellnitz\par
P. Newton
\end{center}

\clearpage
\section*{应用数学教材丛书}
\addcontentsline{toc}{section}{应用数学教材丛书}
\begin{enumerate}
\small
\item Sirovich: 应用数学导论.
\item Wiggins: 应用非线性动力系统与混沌导论.
\item Hale/Koçak: 动力系统与分岔.
\item Chorin/Marsden: 流体力学的数学导论, 第三版.
\item Hubbard/West: 微分方程: 动力系统方法——常微分方程.
\item Sontag: 数学控制理论: 确定性有限维系统, 第二版.
\item Perko: 微分方程与动力系统, 第三版.
\item Seaborn: 超几何函数及其应用.
\item Pipkin: 积分方程教程.
\item Hoppensteadt/Peskin: 医学与生命科学中的建模与模拟, 第二版.
\item Braun: 微分方程及其应用, 第四版.
\item Stoer/Bulirsch: 数值分析导论, 第三版.
\item Renardy/Rogers: 偏微分方程导论.
\item Banks: 生长与扩散现象: 数学框架及应用.
\item Brenner/Scott: 有限元方法的数学理论, 第三版.
\item Van de Velde: 并发科学计算.
\item Marsden/Ratiu: 力学与对称性导论, 第二版.
\item Hubbard/West: 微分方程: 动力系统方法——高维系统.
\item Kaplan/Glass: 理解非线性动力学.
\item Holmes: 摄动方法导论.
\item Curtain/Zwart: 无限维线性系统理论导论.
\item Thomas: 数值偏微分方程: 有限差分方法.
\item Taylor: 偏微分方程: 基础理论.
\item Merkin: 运动稳定性理论导论.
\item Naber: 拓扑, 几何与规范场: 基础.
\item Polderman/Willems: 数学系统理论导论: 行为方法.
\item Reddy: 泛函分析初步及其在边值问题和有限元中的应用.
\item Gustafson/Wilcox: 高等工程数学的解析与计算方法.
\item Tveito/Winther: 偏微分方程导论: 计算方法.
\end{enumerate}
\noindent(续表见索引之后. )

\clearpage
\begin{titlepage}
\centering
\vspace*{2cm}
{\Large Susanne C. Brenner\qquad L. Ridgway Scott\par}
\vspace{2cm}
{\huge\bfseries 有限元方法的数学理论\par}
\vspace{0.8cm}
{\Large 第三版\par}
\vfill
{\Large Springer\par}
\end{titlepage}

\clearpage
\section*{出版信息}
\begin{minipage}[t]{0.47\textwidth}
\textbf{Susanne C. Brenner}\par
Department of Mathematics and Center for Computation and Technology\par
Louisiana State University\par
Baton Rouge, LA 70803\par
USA\par
\texttt{brenner@math.lsu.edu}
\end{minipage}\hfill
\begin{minipage}[t]{0.47\textwidth}
\textbf{L. Ridgway Scott}\par
University of Chicago\par
Chicago, IL 60637\par
USA\par
\texttt{ridg@uchicago.edu}
\end{minipage}

\vspace{1cm}
\textbf{丛书编辑}\par
J.E. Marsden, California Institute of Technology; L. Sirovich, Brown University; S.S. Antman, University of Maryland.

\vspace{1cm}
ISBN 978-0-387-75933-3\qquad e-ISBN 978-0-387-75934-0\par
DOI: 10.1007/978-0-387-75934-0\par
Library of Congress Control Number: 2007939977\par
Mathematics Subject Classification (2000): 65N30, 65--01, 46N40, 65M60, 74S05.

\vspace{0.8cm}
\copyright\ 2008 Springer Science+Business Media, LLC. 保留所有权利. 未经出版者 Springer Science+Business Media, LLC(233 Spring Street, New York, NY 10013, USA)书面许可, 本作品不得全部或部分翻译或复制, 与评论或学术分析有关的简短摘录除外. 禁止以任何现有或今后开发的信息存储和检索形式, 电子改编, 计算机软件或类似或不同的方法使用本作品.

本出版物中商号, 商标, 服务标志及类似术语的使用, 即使未明确标示, 也不应被视为对其是否受专有权保护表达任何意见.

\vfill
无酸纸印刷\par
9 8 7 6 5 4 3 2 1\par
springer.com

\clearpage
\section*{丛书前言}
\label{front:series-preface}
\addcontentsline{toc}{section}{丛书前言}

数学正在物理科学和生物科学中发挥日益重要的作用, 这使各科学学科之间的边界逐渐模糊, 并使人们重新关注应用数学的现代技术和经典技术. 研究和教学中兴趣的复兴促成了``应用数学教材''(TAM)丛书的创立.

研究前沿的高度活跃自然带来了新课程的发展: 数值与符号计算机系统, 动力系统和混沌等较新技术, 正与应用数学的传统方法相互融合并彼此加强. 因此, 本教材丛书旨在满足这些进展当前和未来的需要, 并鼓励开设新课程.

TAM 将出版适合高年级本科课程和研究生入门课程使用的教材. 它将与 AMS 丛书``应用数学科学''相互补充; 后者侧重高级教材和研究级专著.

\begin{flushright}
Pasadena, California\qquad J.E. Marsden\par
Providence, Rhode Island\qquad L. Sirovich\par
College Park, Maryland\qquad S.S. Antman
\end{flushright}

\clearpage
\section*{第三版前言}
\label{front:preface-third-edition}
\addcontentsline{toc}{section}{第三版前言}

本版新增四节, 主题分别为: BDDC 区域分解预条件子(第~\ref{sec:ch07-bddc-preconditioner} 节), 一个收敛的自适应算法(第~\ref{sec:ch09-convergent-adaptive-algorithm} 节), 内罚方法(第~\ref{sec:ch10-discontinuous} 节), 以及分片 $W_p^1$ 函数的 Poincaré--Friedrichs 不等式(第~\ref{sec:ch10-piecewise-poincare-friedrichs} 节). 我们改进了全书许多地方, 其中不少改进由同事建议, 谨致谢意. 本版还增加了新习题, 并扩充和更新了参考文献表.

部分新增材料源自我们的研究, 感谢 National Science Foundation 的资助. 第一作者还感谢 Alexander von Humboldt Foundation 资助她于 2007 年夏季访问德国, 本版工作正是在此期间完成的. 第二作者还感谢 Université Pierre et Marie Curie 资助他过去数年间访问 Paris, 与本版有关的工作在这些访问期间开展.

第一版前言中曾概述本书用于课程的不同方式. 由于部分章号已有变化, 这里重新表述这些建议. 第~\ref{chap:ch00-basic-concepts}--\ref{chap:ch05-n-dimensional-variational-problems} 章构成一门课程的基本材料(这些章号未改变). 第~\ref{chap:ch06-finite-element-multigrid-methods} 章和第~\ref{chap:ch07-additive-schwarz-preconditioners} 章介绍有限元方程线性系统的高效迭代求解器, 但不包含后续章节必需的材料. 强调算法方面的课程应纳入这两章. 同样, 第~\ref{chap:ch08-max-norm-estimates} 章和第~\ref{chap:ch09-adaptive-meshes} 章也不是后续章节所必需的; 涉及有挑战性分析问题的课程可讲授这两章. 前者发展最大范数误差估计并把它用于非线性问题, 后者引入网格自适应性的概念. 不过, 第~\ref{chap:ch10-variational-crimes} 章对后续章节具有关键作用. 也可以只讲授该章第一节和第三节, 然后进入第~\ref{chap:ch11-planar-elasticity} 章或第~\ref{chap:ch12-mixed-methods} 章, 研究应用中常见的典型微分方程组. 第~\ref{chap:ch13-iterative-techniques-mixed-methods} 章实质上是第~\ref{chap:ch12-mixed-methods} 章的延续. 第~\ref{chap:ch10-variational-crimes}--\ref{chap:ch13-iterative-techniques-mixed-methods} 章构成一门强调力学基本模型课程的核心. 第~\ref{chap:ch14-operator-interpolation-applications} 章是一个稍高层次的独立主题, 只依赖第~\ref{chap:ch00-basic-concepts}--\ref{chap:ch05-n-dimensional-variational-problems} 章. 它发展若干泛函分析技术及其在有限元方法中的应用.

\begin{flushright}
Baton Rouge, LA\qquad Susanne C. Brenner\par
Chicago, IL\qquad L. Ridgway Scott\par
20/07/2007
\end{flushright}

\clearpage
\section*{第二版前言}
\label{front:preface-second-edition}
\addcontentsline{toc}{section}{第二版前言}

本版新增两章. 第一章讨论加性 Schwarz 理论及其在多层预条件子和区域分解预条件子中的应用, 第二章介绍后验误差估计量和自适应性. 我们还新增了一节一维自适应网格的例子, 一节离散 Sobolev 不等式, 并在全书增加了新习题. 参考文献表也得到扩充和更新.

我们借此机会再次感谢多年来为本书提供评论和建议的每一个人, 并感谢 National Science Foundation 的资助. 还要感谢 Achi Dosanjh 以及 Springer-Verlag 的制作人员所表现的耐心和细致.

\begin{flushright}
Columbia, SC\qquad Susanne C. Brenner\par
Chicago, IL\qquad L. Ridgway Scott\par
20/02/2002
\end{flushright}

\clearpage
\section*{第一版前言}
\label{front:preface-first-edition}
\addcontentsline{toc}{section}{第一版前言}

本书发展有限元方法的基本数学理论. 有限元方法是工程设计与分析中使用最广泛的技术. 本书的一个目的是把本领域研究者普遍使用但从未发表的基本工具形式化. 本书主要面向数学研究生以及具有较高数学素养的工程师和科学家.

本书曾作为 The University of Michigan, Penn State University 和 University of Houston 的研究生课程教材. 先修要求仅为一门实变量课程, 而对准备充分的工程师和科学家, 在许多情形下甚至不需要这一先修课. 本书可用于一门介绍基本泛函分析, 逼近论和数值分析的课程, 同时以实变量理论的基本技术为基础并加以应用.

第~\ref{chap:ch00-basic-concepts}--\ref{chap:ch05-n-dimensional-variational-problems} 章构成课程的基本材料. 第~\ref{chap:ch00-basic-concepts} 章在一维情形中给出后续内容的缩影. 第~\ref{chap:ch01-sobolev-spaces}--\ref{chap:ch04-polynomial-approximation} 章给出基本理论, 第~\ref{chap:ch05-n-dimensional-variational-problems} 章发展该理论的基本应用. 从这里起, 课程可以向不同方向分流. 第~\ref{chap:ch06-finite-element-multigrid-methods} 章介绍有限元方程线性系统的高效迭代求解器. 虽然从实践角度看这是不可缺少的(这也是我们把它放在显著位置的原因), 但可以跳过, 因为后续章节并不依赖它. 同样, 原第一版第 7 章推导最大范数误差估计并说明这类估计如何用于非线性问题, 也可按需跳过.

不过, 原第一版第 8 章对随后各章具有关键作用. 也可以只讲授该章第一节和第三节, 然后进入原第一版第 9 章, 以考察应用中常见的更复杂微分方程组. 原第一版第 10 章在一定程度上依赖第 9 章, 第 11 章实质上是第 10 章的延续. 原第一版第 12 章介绍 Banach 空间插值技术及其在有限元方法收敛结果中的应用. 这是一个稍高层次的独立主题.

更具体地说, 可以选择三条课程路径. 每条路径的第一步都是学习第~\ref{chap:ch00-basic-concepts}--\ref{chap:ch05-n-dimensional-variational-problems} 章. 希望介绍本学科某些``困难估计''的人, 随后可以从原第一版第 6--8 章和第 12 章中选择. 更关注物理应用的人, 可以选择原第一版第 8.1 节, 第 8.3 节和第 9--11 章. 主要关注算法效率和代码开发问题的人, 可以依次学习原第一版第 6, 8, 10 和 11 章.

本书遗漏的内容极多, 以至难以开始列举. 我们试图列出其中最显著且已有优秀著作可供补充的部分.

我们几乎完全避开了含时问题, 部分原因是已有 Thomée (1984) 一书. 我们对不同类型有限元及其相应逼近理论的广泛论述, 与 Thomée 的方法相互补充. 同样, 我们对物理应用的论述主要限于连续介质力学中的线性系统. Johnson (1987) 一书中可找到更为实质性的物理应用.

本书很少讨论自适应性. 这一活跃研究领域见多种会议论文集(参见 Babuška, Chandra 与 Flaherty 1983, 以及 Babuška, Zienkiewicz, Gago 与 de A. Oliveira 1986).

我们强调可采用的离散化形式的多样性(即不同的``有限元''), 并在可能时把它们表示为依赖某一参数(通常为逼近次数)的族. 这样便形成了``高阶''逼近的精神, 尽管我们不把提高逼近次数(如所谓 P-方法和谱元方法中那样)作为获得收敛的手段, 而是把网格细分作为收敛参数. 关于这一方向的其他方案, 可参见 Szabo 与 Babuška (1991) 的近著.

尽管我们简要介绍了混合方法, 但本书并未恰当地反映这一主题的重要性. 关于该主题的全面论述, 可参见 Brezzi 与 Fortin (1991) 的近著.

我们广泛借鉴 Ciarlet (1978) 一书, 既沿用了其中许多思想, 也把它作为各主题进一步发展的参考. 该书最近在 Ciarlet 与 Lions (1991) 中得到更新, 后者还包含一篇关于混合方法的出色综述. 此外, 可以期待该参考文献所属的 Handbook 丛书今后提供有价值的参考材料.

我们借此机会感谢许多人在本书准备的不同阶段以多种方式提供帮助. 许多曾艰难研读本书早期稿本的学生作出了不可估量的贡献. 读过初稿的许多读者会发现他们的具体建议已被采纳.

本书由作者使用 Springer-Verlag 的 \TeX\ 宏包排版处理.

\clearpage
\tableofcontents
\clearpage
\pagenumbering{arabic}
